\documentclass[12pt]{article}
\usepackage[utf8]{inputenc}
\usepackage[T1]{fontenc}
\usepackage[margin=1in]{geometry}
\usepackage{setspace}
\usepackage{xcolor}
\usepackage[colorlinks=true, urlcolor=blue!60!black, linkcolor=blue!60!black]{hyperref}
\setlength{\parindent}{2em}
\setlength{\parskip}{1.25em}
\renewcommand{\baselinestretch}{1.0}
\begin{document}
\begin{center}
{\LARGE \textit{Chiles v. Salazar}, 607 U.S. \_\_\_ (2026)\par}
\vspace{0.8em}
{\large Case Brief --- Four-Page Edition\par}
\end{center}
\vspace{1em}

\section{Memory Jogger}

The Supreme Court held 8--1 that Colorado's Minor Conversion Therapy Law, as applied to a licensed counselor's pure talk therapy, discriminates based on viewpoint and so triggers strict scrutiny rather than rational-basis review as a regulation of professional conduct (\href{https://www.supremecourt.gov/opinions/25pdf/24-539_fd9g.pdf}{Slip Op.}).

\section{Detailed Case Facts}

Kaley Chiles, a licensed Colorado counselor, provides exclusively talk therapy, with no medication, devices, or aversive techniques. Some minor clients come to her hoping to ``reduce or eliminate unwanted sexual attractions, change sexual behaviors, or grow in the experience of harmony with [their] bod[ies].'' In 2019 Colorado enacted the Minor Conversion Therapy Law (MCTL), which bars licensed counselors from ``any practice or treatment \ldots\ that attempts \ldots\ to change an individual's sexual orientation or gender identity,'' while permitting ``[a]cceptance, support, and understanding'' of identity exploration and ``[a]ssistance to a person undergoing gender transition.'' Violations risk fines, probation, or license revocation. Chiles brought a pre-enforcement, as-applied First Amendment challenge limited to her voluntary talk therapy, conceding the State's power to ban aversive physical interventions. The circuits were split over whether such bans regulate conduct (rational basis) or speech (heightened scrutiny) after \textit{NIFLA v. Becerra} (2018).

\section{Procedural History}

The District of Colorado found standing but denied a preliminary injunction, treating the MCTL as a regulation of professional conduct subject to rational-basis review. The Tenth Circuit affirmed, 2--1, 116 F.4th 1178 (2024); Judge Hartz dissented that recasting talk therapy as conduct was ``a labeling game.'' The Supreme Court granted certiorari, heard argument October 7, 2025, and reversed on March 31, 2026.

\section{Judicial Votes}

Reversed and remanded, \textbf{8--1}.
\begin{itemize}
\item \textbf{Majority (8):} Justice Gorsuch, joined by Chief Justice Roberts and Justices Thomas, Alito, Sotomayor, Kagan, Kavanaugh, and Barrett.
\item \textbf{Concurrence:} Justice Kagan, joined by Justice Sotomayor.
\item \textbf{Dissent:} Justice Jackson.
\end{itemize}

\section{Holding}

A licensing law barring counselors from conversion therapy, as applied to pure talk therapy, regulates speech based on content and viewpoint---not conduct with an incidental burden on speech---and so triggers strict scrutiny. States cannot evade that scrutiny by relabeling speech as a ``treatment'' or ``therapeutic modality,'' or by aggregating licensing, informed-consent, and malpractice traditions into an exemption for ``substandard care.'' The Court found standing but did not itself apply strict scrutiny, remanding instead; it left open aversive physical therapies and viewpoint-neutral regulation of healthcare speech.

\section{Analysis of Opinions}

\textbf{Majority (Gorsuch).} Content-based restrictions trigger strict scrutiny under \textit{Reed}, and viewpoint discrimination is an ``egregious form'' of content regulation. Because ``all Ms. Chiles does is speak---and, as far as she is concerned, speech is all Colorado seeks to regulate,'' the MCTL lets her affirm a minor's identity but forbids helping a consenting minor pursue change. The speech-incident-to-conduct doctrine did not apply: life coaches could give the same advice, and the law targeted Chiles's viewpoint. Nor was there a historic tradition exempting ``substandard care'': counselor licensure dates only to 1976, informed consent attaches to separate physical procedures, and malpractice demands proof of causation and injury, giving ``breathing room for protected speech.'' Citing the APA's 1952 classification of homosexuality and \textit{Buck v. Bell}, Gorsuch warned that professional consensus cannot become enforced orthodoxy.

\textbf{Kagan, with Sotomayor (concurring).} The MCTL was ``textbook'' viewpoint discrimination; a mirror-image law banning affirming therapy would fare the same. She reserved how courts should treat content-based but viewpoint-neutral healthcare rules, which may escape strict scrutiny.

\textbf{Jackson (dissenting).} Starting from \textit{Lambert v. Yellowley}, she stressed the medical consensus that conversion therapy harms minors and read \textit{NIFLA} and \textit{Casey} to turn on whether the State regulates ``speech as speech'' or regulates treatment within a provider-patient relationship. In clinical care ``we insist upon competence, not debate.'' The ruling ``plays with fire'' by making speech-only therapies unregulatable.

\section{Examples: Future Applications}

\begin{enumerate}
\item \textbf{Same-Side: Ban on Affirming Talk Therapy.} A State bars therapy affirming a minor's transgender identity while permitting counseling toward alignment with biological sex. Strict scrutiny applies; both majority and concurrence say the mirror-image law is equally viewpoint-based.
\item \textbf{Same-Side: Counseling on Abortion or Aid in Dying.} A State forbids counselors to speak favorably of these options while allowing the opposite view. Spoken counseling favoring one viewpoint triggers strict scrutiny.
\item \textbf{Opposite-Side: Aversive Physical Techniques.} A ban on electric shocks or nausea-inducing drugs to change orientation regulates non-expressive conduct and survives rational-basis review, as Chiles conceded.
\item \textbf{Opposite-Side: Informed Consent Before Surgery.} Requiring surgeons to disclose factual, uncontroversial risks remains exempt, because it is tied to a separate physical procedure.
\item \textbf{Fence-Sitter: Viewpoint-Neutral ``Predetermined-Outcome'' Ban.} A State bars directing a patient toward any predetermined orientation or identity outcome, or disciplines ``incompetent'' talk therapy. Under the Kagan concurrence it might escape strict scrutiny; under the Jackson dissent such rules inevitably judge viewpoint in application.
\end{enumerate}

\section{Critique}

\textbf{Progressive and health-law views.} Echoing Justice Jackson and Craig Konnoth, critics call the speech-versus-procedure line clinically incoherent, since much of medicine is spoken diagnosis and advice. They question the majority's malpractice distinction, since Colorado applies a preponderance standard to both, and note the irony of invoking \textit{Buck v. Bell} to protect a practice every major medical association condemns.

\textbf{Originalist and free-speech views.} Free-speech advocates praise \textit{Chiles} as a principled application of \textit{NIFLA}, \textit{Reed}, and \textit{Rosenberger} that denies legislatures a blank check to enforce orthodoxy through ``medical consensus.'' Some originalists see unresolved tension in its narrow historical framing and how it fits traditional licensing rules against dangerous verbal therapies.

\section{Key Quotations}

\begin{quote}
``Her speech does not become conduct just because the State may call it that. Nor does her speech become conduct just because it can also be described as a `treatment,' a `therapeutic modality,' or anything else. The First Amendment is no word game. And the rights it protects cannot be renamed away or their protections nullified by `mere labels.' ''
\end{quote}
\noindent --- Justice Gorsuch, Opinion of the Court, slip op. at 12--13.

\begin{quote}
``Medical consensus, too, is not static; it evolves and always has. A prevailing standard of care may reflect what most practitioners believe today, but it cannot mark the outer boundary of what they may say tomorrow.''
\end{quote}
\noindent --- Justice Gorsuch, Opinion of the Court, slip op. at 22.

\begin{quote}
``The Constitution does not pose a barrier to reasonable regulation of harmful medical treatments just because substandard care comes via speech instead of scalpel.''
\end{quote}
\noindent --- Justice Jackson, Dissenting, slip op. at 35.

\end{document}
